% ============================================================
% CAPÍTULO 5: PROTEÍNAS ESTRUCTURALES Y CATALÍTICAS
% ============================================================
% CONTENIDO BASADO EN: Unidad de Aprendizaje II, Temas "Aminoácidos, péptidos y proteínas" y "Proteínas estructurales y catalíticas"
% PROPÓSITO: Distinguir las conformaciones de las proteínas, sus funciones, y verificar el proceso de desnaturalización/renaturalización y los factores que afectan la actividad enzimática.
% ============================================================

\section{Estructura de las proteínas}

% --- CONSEJO: Amplía el concepto de estructura proteica. Diferencia entre proteínas globulares y fibrosas. ---
La función de una proteína está determinada por su estructura tridimensional. Esta estructura se describe en cuatro niveles de organización:

\subsection{Estructura primaria}
% --- CONSEJO: Es la secuencia lineal de aminoácidos. ---
\resumebox{definicion}{Estructura primaria de una proteína}{Es la secuencia lineal de aminoácidos que conforman la proteína, unidos por enlaces peptídicos.}

\subsection{Estructura secundaria}
% --- CONSEJO: Son los plegamientos locales de la cadena polipeptídica, como la hélice $\alpha$ y la hoja $\beta$, estabilizados por puentes de hidrógeno. ---
\resumebox{definicion}{Estructura secundaria de una proteína}{Es el plegamiento regular y repetitivo de la cadena polipeptídica, estabilizado por puentes de hidrógeno entre los átomos del esqueleto peptídico. Los dos tipos principales son la hélice $\alpha$ y la hoja $\beta$ plegada.}

\subsection{Estructura terciaria}
% --- CONSEJO: Es la conformación tridimensional completa de una sola cadena polipeptídica, estabilizada por interacciones entre cadenas laterales de aminoácidos. ---
\resumebox{definicion}{Estructura terciaria de una proteína}{Es la conformación tridimensional completa de una cadena polipeptídica, estabilizada por interacciones entre las cadenas laterales de los aminoácidos (fuerzas de Van der Waals, interacciones hidrofóbicas, puentes de hidrógeno, enlaces iónicos y puentes disulfuro).}

\subsection{Estructura cuaternaria}
% --- CONSEJO: Se aplica solo a proteínas con dos o más cadenas polipeptídicas (subunidades). ---
\resumebox{definicion}{Estructura cuaternaria de una proteína}{Es la organización de dos o más subunidades polipeptídicas para formar un complejo proteico funcional.}

\section{Proteínas fibrosas vs. Globulares}

% --- CONSEJO: Haz una comparación clara entre estos dos tipos de proteínas, ya que el programa lo pide explícitamente. ---
\begin{table}[H]
    \centering
    \caption{Comparación entre proteínas fibrosas y globulares}
    \label{tab:proteinas_fibrosas_globulares}
    \begin{tabular}{|p{0.3\linewidth}|p{0.3\linewidth}|p{0.3\linewidth}|}
        \hline
        \theader{Característica} & \theader{Proteínas fibrosas} & \theader{Proteínas globulares} \\
        \hline
        \textbf{Forma} & Alargada, filamentosa o en forma de lámina. & Compacta, esferoidal. \\
        \textbf{Función} & Estructural y de soporte. & Funcional (enzimas, transporte, etc.). \\
        \textbf{Solubilidad} & Insolubles en agua. & Solubles en agua. \\
        \textbf{Estructura} & Principalmente estructura secundaria repetitiva. & Mezcla de estructuras secundarias, con estructura terciaria definida. \\
        \textbf{Ejemplos} & Colágeno, queratina, elastina. & Enzimas, hemoglobina, anticuerpos. \\
        \hline
    \end{tabular}
\end{table}

\section{Enzimas: Proteínas catalíticas}

% --- CONSEJO: Introduce el concepto de enzima, su importancia como catalizadores biológicos y su clasificación. ---
\resumebox{definicion}{Enzimas}{Las enzimas son proteínas que actúan como catalizadores biológicos, acelerando la velocidad de las reacciones químicas en los seres vivos sin consumirse en el proceso.}

\subsection{Clasificación de las enzimas}
% --- CONSEJO: Explica la clasificación de la IUBMB (Unión Internacional de Bioquímica y Biología Molecular) y da ejemplos de cada clase. ---
\begin{enumerate}
    \item \textbf{Oxidorreductasas:} Catalizan reacciones de oxidación-reducción (transferencia de electrones). Ej. Deshidrogenasas, oxidasas.
    \item \textbf{Transferasas:} Catalizan la transferencia de grupos funcionales. Ej. Quinasas, transaminasas.
    \item \textbf{Hidrolasas:} Catalizan la hidrólisis de enlaces (ruptura con agua). Ej. Proteasas, lipasas.
    \item \textbf{Liasas:} Catalizan la adición o eliminación de grupos sin hidrólisis ni oxidación. Ej. Descarboxilasas.
    \item \textbf{Isomerasas:} Catalizan la isomerización de moléculas. Ej. Fosfoglucosa isomerasa.
    \item \textbf{Ligasas:} Catalizan la unión de dos moléculas con hidrólisis de ATP. Ej. Sintetasas.
\end{enumerate}

\section{Factores que afectan la actividad enzimática}

% --- CONSEJO: Explica cómo el pH, la temperatura y la concentración de sustrato afectan la velocidad de una reacción enzimática. ---
\begin{itemize}
    \item \textbf{pH y temperatura:} Las enzimas tienen un pH y una temperatura óptimos donde su actividad es máxima. Fuera de estos rangos, la enzima se desnaturaliza.
    \item \textbf{Concentración de sustrato:} Al aumentar la concentración de sustrato, la velocidad de la reacción aumenta hasta alcanzar una velocidad máxima (Vmax), donde la enzima está saturada.
    \item \textbf{Inhibidores:} Moléculas que se unen a la enzima y disminuyen su actividad. Pueden ser competitivos, no competitivos o acompetitivos.
    \item \textbf{Cofactores y coenzimas:} Muchas enzimas necesitan iones metálicos (cofactores) o moléculas orgánicas (coenzimas) para ser activas.
\end{itemize}

\section{Desnaturalización y renaturalización de proteínas y enzimas}

% --- CONSEJO: Explica el concepto de desnaturalización como la pérdida de la estructura tridimensional, y la renaturalización como el replegamiento correcto. ---
\resumebox{definicion}{Desnaturalización de proteínas}{Es la pérdida de la estructura tridimensional (secundaria, terciaria y/o cuaternaria) de una proteína, manteniendo su estructura primaria intacta. Generalmente es un proceso irreversible.}

La \textbf{renaturalización} (o replegamiento) es el proceso por el cual una proteína desnaturalizada recupera su conformación nativa y, por lo tanto, su actividad biológica. No todas las proteínas pueden renaturalizarse espontáneamente; algunas requieren la ayuda de proteínas llamadas chaperonas.

% --- CONSEJO: Aquí puedes incluir un ejemplo clásico como la ribonucleasa A, que puede desnaturalizarse y renaturalizarse in vitro. ---

\section{Funciones e importancia biotecnológica de las proteínas estructurales y catalíticas}

% --- CONSEJO: Amplía la importancia de las proteínas más allá de lo básico, enfocándote en su aplicación biotecnológica. ---
\begin{itemize}
    \item \textbf{Proteínas estructurales:} Se utilizan en la industria textil (queratina, colágeno), en la producción de biomateriales (colágeno para implantes) y en la ingeniería de tejidos.
    \item \textbf{Enzimas:} Son la base de muchos bioprocesos:
    \begin{itemize}
        \item \textbf{Industria alimentaria:} Uso de amilasas, proteasas y lipasas.
        \item \textbf{Detergentes:} Proteasas y lipasas en detergentes enzimáticos.
        \item \textbf{Diagnóstico:} Enzimas usadas en kits de diagnóstico y biosensores.
        \item \textbf{Biocatálisis:} Uso de enzimas para la síntesis de compuestos orgánicos.
    \end{itemize}
\end{itemize}